\documentclass[11pt]{article}
\usepackage[a4paper,margin=2.5cm]{geometry}
\usepackage[utf8]{inputenc}
\usepackage[T1]{fontenc}
\usepackage{lmodern}
\usepackage{amsmath,amssymb}
\usepackage{booktabs}
\usepackage{xcolor}
\usepackage[numbers]{natbib}
\usepackage[hidelinks]{hyperref}

% Result numbers are filled from the published snapshot (ledger entry and
% docs/validation.md). Until then every number is a visible placeholder.
\newcommand{\TBD}[1]{\textcolor{red}{\textbf{[#1]}}}
% Result values. Validation numbers come from docs/validation.md (kairos run
% validation, 2026-11-07); dial numbers come from build/snapshot.json and
% ledger entry 0, recorded by the weekly GitHub Actions run on 2026-11-07
% (commit 8a383b9). Do not edit numbers anywhere else in the paper.
\newcommand{\resAsOf}{1~September 2026}
% population: t_c_1960 = 2052.210, 95% interval (2018.980, 2113.880).
\newcommand{\resPopTc}{2052.2 (95\% interval 2019.0--2113.9)}
% population: local delta 3.737 in 1960, -0.685 in 1980 (trend -0.168/yr).
\newcommand{\resPopBendYear}{1980}
% dotcom: t_c_last_pre_peak = 2000.423; published window 2000.0--2000.5
% (Johansen & Sornette 2000).
\newcommand{\resDotcomTc}{2000.42 vs.\ 2000.0--2000.5}
% bitcoin: lead_time_days = 350 (upper bound: baseline LPPLS alarms persistently).
\newcommand{\resBtcLead}{350 days}
% n/a: the COVID case scores mean absolute log error, not CRPS. Lockdown-
% conditioned MAE(log) 1.062 vs. no-event counterfactual 14.78.
\newcommand{\resCovidCrps}{n/a}
% moore: 0 alarms over 37 yearly as-of dates.
\newcommand{\resMooreFAR}{0/37}
% ai_history: latest_breaks = (2017.080, 2018.580).
\newcommand{\resAiChangepoints}{2017.1, 2018.6}
% n/a: no validation case computes interval coverage (scoring.interval_coverage
% exists but the runner does not call it).
\newcommand{\resCoverage}{n/a}
\newcommand{\resDelta}{0.081 (90\% CI 0.026--0.085)}
\newcommand{\resKappa}{0.369 (90\% CI 0.111--0.549)}
\newcommand{\resHstar}{23.0 months (90\% CI 22.2--23.8)}
\newcommand{\resGapVelocity}{0.443 per year}
\newcommand{\resLevel}{0}
% n/a: snapshot alert.lead_time is null at level 0 (no lead time is issued).
\newcommand{\resLeadTime}{n/a}
\newcommand{\resLedgerHash}{9e3b99b1c60116bcb35812425bebc84d6d30d5a9e81683776e4b0a0b5d31ddcb}


\title{Measuring the Distance to a Singularity:\\
Forecastability as an Early-Warning Signal}
\author{KAIROS Observatory contributors\\
\small\url{https://github.com/guru-bharadwaj20/Kairos-Observatory}}
\date{Draft, \today}

\begin{document}
\maketitle

\begin{abstract}
A finite-time singularity is a time-series event: growth that is faster than
exponential reaches infinity at a finite critical date $t_c$. Vinge's original
definition of the technological singularity is also about time series: it is the
point beyond which the future can no longer be forecast. We combine both
readings in KAIROS, an open instrument that tracks AI capability $C(t)$ against
human oversight capacity $H(t)$. It reports four dials: an acceleration
exponent $\delta$ with critical time $t_c$; a recursion coefficient $\kappa$;
the Vinge Horizon $h^*$, the longest horizon still forecastable within a fixed
factor error; and the velocity of the Jonas Gap $G=\ln C-\ln H$. Decision rules
with published parameters turn the dials into pre-agreed tripwire levels, and
every weekly forecast is hashed into a timestamped public ledger before its
outcome is known. In backtests the method recovers von Foerster's 1960
doomsday date (\resPopTc), sees $\delta$ fall toward zero by \resPopBendYear,
and stays at level~0 on Moore's law (false-alarm rate \resMooreFAR). The
current reading is level \resLevel{} with a lead time of \resLeadTime.
\end{abstract}

\section{Introduction}

In 1960 von Foerster, Mora and Amiot fitted a hyperbola to world population and
found that it diverged on Friday, 13 November 2026~\cite{vonfoerster1960}.
The singularity was real in the data and was avoided in reality: the curve bent
during the demographic transition. Two lessons follow. Finite-time
singularities can be seen in data decades ahead, and they can be bent. The hard
part is telling early whether a curve is bending.

AI governance faces the Collingridge dilemma~\cite{collingridge1980}: control is
easy while harms are unknown and hard once they are clear. Task horizons of AI
agents have doubled roughly every seven months~\cite{metr2025}, and frontier
training compute grows several-fold per year~\cite{epoch}. What is missing is a
neutral instrument that issues falsifiable warnings tied to responses agreed in
advance. This paper describes one and reports its validation.

Our contributions are: (i)~a forecastability-based definition of distance to a
singularity, the Vinge Horizon $h^*$; (ii)~a relational definition of danger,
the Jonas Gap velocity $dG/dt$; (iii)~decision rules written as mathematics with
published parameters; and (iv)~a pre-registration ledger that makes a
self-defeating prophecy scorable through conditional forecasts.

\section{Method}

\subsection{Null model}
Prophet~\cite{taylor2018} is fitted to $\ln y$. A piecewise-linear trend on the
log scale is piecewise exponential growth. It cannot represent a finite-time
singularity, and that is the point: it stands for the boring world, and
persistent one-sided residuals against it are evidence of acceleration.

\subsection{Acceleration: $\delta$ and $t_c$}
Hyperbolic growth $\dot x = k x^{1+\delta}$ gives
$x(t)=x_0(1-t/t_c)^{-1/\delta}$. With $g=d\ln x/dt$,
\begin{equation}
\ln g = \ln k + \delta \ln x, \qquad t_c - t = \frac{1}{\delta\, g}.
\end{equation}
$g$ is estimated from observed data with a cross-validated smoother, not from a
forecast trend, whose slope is piecewise constant by construction and biases
$\delta$ to zero. Because $\ln g$ and $\ln x$ share autocorrelated errors,
intervals come from a moving block bootstrap that refits the smoother on every
resample. A direct nonlinear fit of $x(t)$ and the log-periodic power law
singularity (LPPLS) model~\cite{sornette2003}
\begin{equation}
\ln x = A + B(t_c-t)^m + C(t_c-t)^m\cos\!\big(\omega\ln(t_c-t)-\phi\big)
\end{equation}
give independent estimates of $t_c$. Significance across metrics is controlled
with Benjamini--Hochberg, and metrics count as independent only if they come
from different source groups.

\subsection{Recursion: $\kappa$}
NeuralProphet~\cite{triebe2021} is fitted with exogenous drivers only (compute,
investment, talent), then with capability's own lags added. $\kappa$ is the
relative out-of-sample skill gain, reported against a permutation null with
shuffled self-lags and cross-checked with Granger tests.

\subsection{The Vinge Horizon $h^*$}
Rolling-origin backtests with data vintages give forecast error as a function of
horizon. Error is measured in log units, as absolute log error and CRPS, with a
tolerance stated as a factor (for example within $2\times$). $h^*$ is the
\emph{first} horizon at which the smoothed error crosses the tolerance;
taking the maximum passing horizon would let one lucky long horizon inflate it.
$h^*(t)$ is recomputed at every monthly cutoff and tested for shrinkage with
Mann--Kendall and Theil--Sen. This is Lorenz's predictability limit~\cite{lorenz1963}
applied to technology.

\subsection{The Jonas Gap}
$G(t)=\ln C(t)-\ln H(t)$, with both composite indices normalised at a base year.
The level of $G$ is arbitrary; only $dG/dt$ is interpreted. The dial is in danger
when $dG/dt>0$ is sustained under a run-length rule and a Newey--West $t$-test.

\subsection{Decision rules}
A missed catastrophe costs $\lambda$ times a false alarm~\cite{jonas1984}, so a
dial raises its alarm when $P(\text{danger}) > 1/(1+\lambda)$, with
probabilities taken from bootstrap and ensemble distributions. The maximin
decision over the scenario ensemble~\cite{rawls1971} is reported next to the
expected-value one. Lead time is the time until the upper forecast quantile
crosses the danger threshold. Tripwire levels 0--4 count dials in danger
(Table~\ref{tab:levels}). $\lambda$ and every threshold are versioned YAML.

\begin{table}[h]
\centering
\begin{tabular}{@{}lll@{}}
\toprule
Level & Trigger & Recommended response\\
\midrule
0 Normal & All dials calm & Weekly publication\\
1 Watch & One dial & More data, public notice\\
2 Warning & Two dials, or one on $\ge 3$ metrics & Mandatory third-party evaluation\\
3 Alarm & Three dials & Pause scaling past the frontier\\
4 Kairos & All four dials & Last cheap window\\
\bottomrule
\end{tabular}
\caption{Tripwire levels, agreed in advance.}
\label{tab:levels}
\end{table}

\subsection{Pre-registration}
Each weekly forecast is serialised as canonical JSON, hashed with SHA-256,
appended to a hash-chained ledger that includes the data manifest hash, and
anchored with OpenTimestamps~\cite{popper1959}. Forecasts are conditional (``if
no new intervention, $X$ reaches $Y$ by $Z$''), and interventions are logged as
events, so a warning that is acted on can be scored against the counterfactual
instead of counting as a false alarm~\cite{merton1948}.

\section{Validation}

\begin{table}[h]
\centering
\begin{tabular}{@{}lll@{}}
\toprule
Case & Check & Result\\
\midrule
World population & $t_c$ from 1960 data & \resPopTc\\
 & $\delta$ falls toward 0 by & \resPopBendYear\\
Nasdaq 2000 & LPPLS $t_c$ vs.\ Sornette & \resDotcomTc\\
Bitcoin 2017 & Warning lead time & \resBtcLead\\
COVID-19 2020 & Counterfactual CRPS gain & \resCovidCrps\\
Moore's law & False-alarm rate & \resMooreFAR\\
AI history & Changepoints found blind & \resAiChangepoints\\
All & 90\% interval coverage & \resCoverage\\
\bottomrule
\end{tabular}
\caption{Historical validation. Numbers come from \texttt{docs/validation.md}.}
\end{table}

\section{Current reading}
As of \resAsOf, the dials read $\delta=\resDelta$, $\kappa=\resKappa$,
$h^*=\resHstar$ and $dG/dt=\resGapVelocity$. The tripwire level is
\resLevel{} and the lead time is \resLeadTime. The ledger entry is
\texttt{\resLedgerHash}.

\section{Limitations}
Frontier data are sparse, delayed or secret; we use many proxies and treat
growing secrecy as a signal. Metrics can be gamed; held-out metrics rotate and a
divergence between public and held-out trends is flagged. Singularity fits on
noisy data raise false alarms; bootstrap intervals, multi-dial confirmation and
a published false-alarm rate mitigate this. The clock could feed race dynamics,
so $H$ and the recommended responses are always shown with $C$. KAIROS cannot
stop anything; it buys lead time, and the decisions stay with people.

\section{Conclusion}
Forecastability is measurable, and its decline is a model-free warning sign.
Combined with acceleration, recursion and the capability--comprehension gap, and
bound to pre-registered forecasts, it gives an instrument that is built to be
proven wrong.

\bibliographystyle{plainnat}
\bibliography{refs}
\end{document}
